% GENERATED FILE. Edit canonical lessons, then run npm run generate:books.
% canonical-source-hash: fnv1a64:262d976308b64aa0
% canonical-lessons: FR-C396-installer, FR-C396-desinstaller, FR-C396-appareil, FR-C396-logiciel, FR-C396-mise-a-jour

\chapter{Installer un logiciel}
\label{ch:fr-installer-un-logiciel}
\hlchaptermodality{396}

\begin{chapteropening}
\textbf{By the end of this chapter:} \emph{I can install and remove software, talk about the device it runs on and about updates.}

Complete the last lesson of chapter 396: I can install and remove software, talk about the device it runs on and about updates.
\end{chapteropening}

\section[installer]{installer — to install; to set up}

\label{lesson:FR-C396-installer}



\begin{quote}
Before the new one: say the French for to connect, then the French for to access.
\end{quote}

\subsection*{You'll want to know: installer}
\textbf{installer} — \textquotedblleft{}to install; to set up\textquotedblright{}.

\begin{quote}
J'ai installé une nouvelle application sur mon téléphone.
\end{quote}

\begin{quote}
\emph{I installed a new app on my phone.}
\end{quote}

\emph{S'installer} — to settle in, to move in.

\begin{cousinweb}
\emph{installer} is usually traced to Medieval Latin \emph{installāre} \textquotedblleft{}to put into a stall, to install in office\textquotedblright{}, built on a Germanic word for a stall; EtymDB lists relatives such as Old English \emph{steall} \textquotedblleft{}a standing place\textquotedblright{}.
\end{cousinweb}

\begin{grammarlens}[title={What you've built}]
One more word for telling, explaining and describing at B1.
\end{grammarlens}

\subsection*{Your turn}
\begin{itemize}
\raggedright
  \item \emph{Say it:} \emph{installer}
  \item \emph{Say it:} \emph{installer}, once more
  \item \emph{Say it:} a sentence of your own with \emph{installer}
  \item \emph{Recall:} say the French for to connect, then the French for to access, then say \emph{installer} again
\end{itemize}

\begin{tcolorbox}[breakable,colback=teal!4,colframe=teal!35!black,title={Before you move on}]
What does installer mean? (\textquotedblleft{}To install; to set up\textquotedblright{}.) Say it once more.
\end{tcolorbox}
\section[désinstaller]{désinstaller — to uninstall}

\label{lesson:FR-C396-desinstaller}



\begin{quote}
Before the new one: say the French for to access, then the French for to install.
\end{quote}

\subsection*{You'll want to know: désinstaller}
\textbf{désinstaller} — \textquotedblleft{}to uninstall\textquotedblright{}.

\begin{quote}
J'ai désinstallé ce jeu : il prenait trop de place.
\end{quote}

\begin{quote}
\emph{I uninstalled that game: it took up too much space.}
\end{quote}

\emph{Dés-} + \emph{installer}: the opposite.

\begin{cousinweb}
\emph{désinstaller} is built from \emph{dés-} and \emph{installer}, which is usually traced to Medieval Latin \emph{installāre} \textquotedblleft{}to put into a stall\textquotedblright{}, built on a Germanic word for a stall.
\end{cousinweb}

\begin{grammarlens}[title={What you've built}]
One more word for telling, explaining and describing at B1.
\end{grammarlens}

\subsection*{Your turn}
\begin{itemize}
\raggedright
  \item \emph{Say it:} \emph{désinstaller}
  \item \emph{Say it:} \emph{désinstaller}, once more
  \item \emph{Say it:} a sentence of your own with \emph{désinstaller}
  \item \emph{Recall:} say the French for to access, then the French for to install, then say \emph{désinstaller} again
\end{itemize}

\begin{tcolorbox}[breakable,colback=teal!4,colframe=teal!35!black,title={Before you move on}]
What does désinstaller mean? (\textquotedblleft{}To uninstall\textquotedblright{}.) Say it once more.
\end{tcolorbox}
\section[l'appareil]{l'appareil — the device, the appliance}

\label{lesson:FR-C396-appareil}



\begin{quote}
Before the new one: say the French for to install, then the French for to uninstall.
\end{quote}

\subsection*{You'll want to know: l'appareil}
\textbf{l'appareil} — \textquotedblleft{}the device, the appliance\textquotedblright{}.

\begin{quote}
Cet appareil ne marche plus.
\end{quote}

\begin{quote}
\emph{This device does not work any more.}
\end{quote}

\emph{Un appareil photo} — a camera. On the phone, \emph{qui est à l'appareil ?} — who is speaking? Masculine.

\begin{cousinweb}
\emph{appareil} is formed on \emph{appareiller} \textquotedblleft{}to prepare, to fit out\textquotedblright{} (EtymDB), from a Vulgar Latin \emph{appariculāre}, built on Latin \emph{parāre} \textquotedblleft{}to prepare\textquotedblright{} (Lewis \& Short).
\end{cousinweb}

\begin{grammarlens}[title={What you've built}]
One more word for telling, explaining and describing at B1.
\end{grammarlens}

\subsection*{Your turn}
\begin{itemize}
\raggedright
  \item \emph{Say it:} \emph{l'appareil}
  \item \emph{Say it:} \emph{l'appareil}, once more
  \item \emph{Say it:} a sentence of your own with \emph{l'appareil}
  \item \emph{Recall:} say the French for to install, then the French for to uninstall, then say \emph{l'appareil} again
\end{itemize}

\begin{tcolorbox}[breakable,colback=teal!4,colframe=teal!35!black,title={Before you move on}]
What does l'appareil mean? (\textquotedblleft{}The device, the appliance\textquotedblright{}.) Say it once more.
\end{tcolorbox}
\section[le logiciel]{le logiciel — the software, the program}

\label{lesson:FR-C396-logiciel}



\begin{quote}
Before the new one: say the French for to uninstall, then the French for the device.
\end{quote}

\subsection*{You'll want to know: le logiciel}
\textbf{le logiciel} — \textquotedblleft{}the software, the program\textquotedblright{}.

\begin{quote}
Ce logiciel est gratuit.
\end{quote}

\begin{quote}
\emph{This software is free.}
\end{quote}

Québec also says \emph{le programme}.

\begin{cousinweb}
\emph{logiciel} is usually described as a French coinage of the twentieth century from \emph{logique} \textquotedblleft{}logic\textquotedblright{} and \emph{-iel}, made to replace English \emph{software}. \emph{logique} comes from Greek \emph{logik\'{\={e}}} \textquotedblleft{}the art of reasoning\textquotedblright{}.
\end{cousinweb}

\begin{grammarlens}[title={What you've built}]
One more word for telling, explaining and describing at B1.
\end{grammarlens}

\subsection*{Your turn}
\begin{itemize}
\raggedright
  \item \emph{Say it:} \emph{le logiciel}
  \item \emph{Say it:} \emph{le logiciel}, once more
  \item \emph{Say it:} a sentence of your own with \emph{le logiciel}
  \item \emph{Recall:} say the French for to uninstall, then the French for the device, then say \emph{le logiciel} again
\end{itemize}

\begin{tcolorbox}[breakable,colback=teal!4,colframe=teal!35!black,title={Before you move on}]
What does le logiciel mean? (\textquotedblleft{}The software, the program\textquotedblright{}.) Say it once more.
\end{tcolorbox}
\section[la mise à jour]{la mise à jour — the update}

\label{lesson:FR-C396-mise-a-jour}



\begin{quote}
Before the new one: say the French for the device, then the French for the software.
\end{quote}

\subsection*{You'll want to know: la mise à jour}
\textbf{la mise à jour} — \textquotedblleft{}the update\textquotedblright{}.

\begin{quote}
Mon téléphone fait une mise à jour cette nuit.
\end{quote}

\begin{quote}
\emph{My phone is doing an update tonight.}
\end{quote}

\emph{Mettre à jour} — to update. Plural \emph{les mises à jour}.

\begin{cousinweb}
\emph{mise à jour} is a putting up to date. \emph{mise} is the participle of \emph{mettre}, from Latin \emph{mittere} \textquotedblleft{}to send, to put\textquotedblright{}, and \emph{jour} is inherited from Latin \emph{diurnum} \textquotedblleft{}daily\textquotedblright{} (EtymDB).
\end{cousinweb}

\begin{grammarlens}[title={What you've built}]
One more word for telling, explaining and describing at B1.
\end{grammarlens}

\subsection*{Your turn}
\begin{itemize}
\raggedright
  \item \emph{Say it:} \emph{la mise à jour}
  \item \emph{Say it:} \emph{la mise à jour}, once more
  \item \emph{Say it:} a sentence of your own with \emph{la mise à jour}
  \item \emph{Recall:} say the French for the device, then the French for the software, then say \emph{la mise à jour} again
\end{itemize}

\begin{tcolorbox}[breakable,colback=teal!4,colframe=teal!35!black,title={Before you move on}]
What does la mise à jour mean? (\textquotedblleft{}The update\textquotedblright{}.) Say it once more.
\end{tcolorbox}
